\documentclass[beginterm]{../package/riaeproposal}

\thesistilte{这是一个特别特别特别特别特别长长的课题小于二十五字}
\title{\buaaatthesistitle}
\authorname{研究生姓名}
\author{\buaaatauthorname}

\tutor{导师姓名 \quad 职称}
\stuno{XXXX}
\major{专业名称}

\begin{document}
\pdfbookmark[0]{\buaaatthesistitle 开题报告}{cover}
\maketitle

\linespread{1.5}
\pagestyle{frontmatter}
\maketoc 

\section{论文选题依据}
\clearpage
\section{国内外研究现状}
\clearpage
\section{学位论文的研究方案}
\subsection{研究目标}

\subsection{研究内容}
\subsection{解决问题}
\subsection{技术路线}
\clearpage
\section{关键技术及难点}
\clearpage
\section{预期成果}
\clearpage
\section{工作进度安排}
论文工作安排如下：\begin{itemize}[leftmargin=9em]
    \item[2022.10-2022.12] 查阅国内外相关文献，进行论文调研
    \item[2023.01-2023.02] 搭建、调试数值仿真环境，进行可压缩湍流数值直接数值模拟计算的测试
    \item[2023.02-2023.08]  在开展毕业实习同时，进行二维可压缩湍流的直接数值模拟工作，并基于数值仿真结果进行可压缩湍流统计量的相关研究
    \item[2023.08-2023.11]  基于数值仿真结果与统计学结果，开展解析理论推导研究
    \item[2023.11-2023.12] 对研究内容进行总结，完成毕业论文撰写与答辩
\end{itemize}
\clearpage
\section{参考文献}
\vspace{-2em}
\bibliography{refs}
\nocite{*}
\end{document}